\documentclass[10pt,a4paper]{amsart}
\usepackage[margin=19mm]{geometry}
\usepackage{amsmath,amssymb,mathtools}
\usepackage[scr=rsfs,cal=euler]{mathalfa}
\usepackage{xfrac}
\usepackage[unicode,hidelinks,bookmarks=false]{hyperref}
\newcommand{\ie}{i.e.\ }
\newcommand{\cf}{cf.\ }
\newcommand{\resp}{respectively\ }
\newcommand{\Subsection}{\subsection}
\providecommand{\nobreakdash}{\nobreak\hbox{-}}
\setlength{\emergencystretch}{2em}
\multlinegap0pt
\usepackage{bm}
\usepackage{bbm}
\usepackage{dsfont}
\usepackage{xspace}
\usepackage{cancel}
\usepackage{mathtools}
\usepackage{amsthm}
\makeatletter
\@ifundefined{theorem}{\newtheorem{theorem}{Theorem}}{}
\@ifundefined{claim}{\newtheorem{claim}[theorem]{Claim}}{}
\@ifundefined{lemma}{\newtheorem{lemma}[theorem]{Lemma}}{}
\makeatother
\title[Minimizing Submodular Functions over Hierarchical Families]{Minimizing Submodular Functions over Hierarchical Families}
\author{Ryuhei Mizutani}
\date{Source: arXiv:2601.14805v2 (CC BY 4.0)}
\begin{document}
\maketitle

\noindent\emph{Source and licence.} Minimizing Submodular Functions over Hierarchical Families. Version arXiv:2601.14805v2 (CC BY 4.0), distributed under the Creative Commons Attribution 4.0 International licence, as recorded in the arXiv licence metadata for this exact version and shown on the abstract page https://arxiv.org/abs/2601.14805v2. Full rights notice: licenses/arxiv-2601.14805-CC-BY-4.0.txt.\par\smallskip

\noindent\textit{Editorial scope.} Counted unit: the Lemma whose source label is \texttt{lem:key} in the article (source0.tex lines 201--204, source label \texttt{lem:key}), delivered together with the article's own complete original proof of it (source0.tex lines 205--250, one \texttt{proof} environment of 46 lines). No mathematics is written, repaired or re-derived: statement and proof are the article's own text line for line. Required context retained uncounted: none. Every definition, display and label of those lines is retained unchanged. Omitted from this package and not redistributed: 1--200, 251--394. Nothing inside a retained excerpt is omitted or altered: every retained body is delivered exactly as the article's lines are written, so no omission receipt of any kind accompanies this package. Citations: the retained excerpts carry no citation command at all, so no bibliography entry of the article is transcribed and no reference is redistributed. Reference adaptation: none was needed and none was made; every reference inside the delivered unit resolves inside it, so no unresolved reference or citation is produced. Printed slips kept literal: no printed word, symbol, hypothesis or number of the counted statement or of its proof was altered. Layout and class: the article's own class is \texttt{article} with packages including inputenc, geometry, amsmath, amssymb, amsfonts, graphicx, amsthm, enumitem, bm, color, complexity, algorithm, algpseudocode, float; the wrapper is the lane's pinned \texttt{amsart} editorial wrapper at 10pt on A4, which restates the theorem-like environments the retained text needs and adds the article's own macro definitions that the retained text needs (none), with extra wrapper packages (bm, bbm, dsfont, xspace, cancel, mathtools). The delivered numbering is the unit's own. Assets: no retained span contains a figure environment, a graphics-inclusion command, a PDF-page inclusion, a table or a TikZ picture; no image file is copied into this package, the package declares no asset dependency and no image file is redistributed with it. Licence: version arXiv:2601.14805v2 of this article is distributed under the Creative Commons Attribution 4.0 International licence, as recorded in the arXiv licence metadata for this exact version and shown on the abstract page https://arxiv.org/abs/2601.14805v2. A full rights notice is delivered at licenses/arxiv-2601.14805-CC-BY-4.0.txt. Characters: every retained excerpt is pure ASCII, so the delivered engine, XeLaTeX with its default Latin Modern text font, reports no missing character.
\begin{lemma}
\label{lem:key}
For any $X\in \mathcal{F}$, there exists $S\subseteq X$ such that $|S|\le k$ and $Y\in \mathcal{F}$ for all $Y\subseteq X$ with $S\subseteq Y$.
\end{lemma}

\begin{proof}
Take any $X\in \mathcal{F}$.
If $|X|\le k$, then setting $S=X$ satisfies the condition, and we are done.
Consider the case when $|X|\ge k+1$.
Since $2^V\setminus \mathcal{F}$ is a $k$-hierarchical lattice, there exist families $\emptyset=\mathcal{H}_0\subseteq \mathcal{H}_1\subseteq  \cdots \subseteq \mathcal{H}_k=2^V\setminus \mathcal{F}$ satisfying the condition $(*)$.
We prove the following claim by induction:
\begin{claim}
\label{clm:induction}
There exist $k$ distinct elements $v_1,v_2,\ldots,v_k\in X$ such that the following condition holds: for each $i\in [k]$, $\{v_1,v_2,\ldots,v_i\}\not\subseteq Y$ for all $Y\in \mathcal{H}_i\setminus \mathcal{H}_{i-1}$ with $Y\subseteq X$.
\end{claim}
\begin{proof}[Proof of Claim \ref{clm:induction}]
We inductively construct $v_1,v_2,\ldots,v_k$ satisfying the condition.
We first show the existence of $v_1\in X$ satisfying the condition.
Suppose to the contrary that for each $v\in X$, there exists $Y\in \mathcal{H}_1$ such that $Y\subseteq X$ and $v\in Y$.
Then, we have
\begin{align*}
\bigcup_{Y\in \mathcal{H}_1,Y\subseteq X}Y=X.
\end{align*}
Combined with the fact that $\mathcal{H}_1$ is a lattice, this implies $X\in \mathcal{H}_1$, which contradicts $X\in \mathcal{F}$.

We next show the existence of $v_2,\ldots,v_k\in X$ satisfying the condition.
As induction hypothesis, assume that there exist $v_1,v_2,\ldots,v_{j-1}$ satisfying the condition for some $2\le j\le k$.
For convenience, let $U_{j-1}:=\{v_1,v_2,\ldots,v_{j-1}\}$.
We show the existence of $v_j\in X$ satisfying the condition.
Suppose to the contrary that for each $v\in X\setminus U_{j-1}$, there exists $Y\in \mathcal{H}_j\setminus \mathcal{H}_{j-1}$ such that $U_{j-1}\cup \{v\}\subseteq Y\subseteq X$.
Define a subfamily $\mathcal{H}_j'$ of $\mathcal{H}_j\setminus \mathcal{H}_{j-1}$ as follows:
\begin{align*}
\mathcal{H}_j':=\{Y\in \mathcal{H}_j\setminus \mathcal{H}_{j-1}\mid U_{j-1}\subsetneq Y\subseteq X\}.
\end{align*}
By the assumption, for each $v\in X\setminus U_{j-1}$, there exists $Y\in \mathcal{H}_j'$ with $v\in Y$.
This implies
\begin{align*}
\bigcup_{Y\in \mathcal{H}_j'}Y=X.
\end{align*}
If $Y_1\cup Y_2\in \mathcal{H}_j\setminus \mathcal{H}_{j-1}$ for each pair of $Y_1,Y_2\in \mathcal{H}_j'$, then by the above equation we have $X\in \mathcal{H}_j\setminus \mathcal{H}_{j-1}$, a contradiction.
Hence, there exist $Y_1,Y_2\in \mathcal{H}_j'$ such that $Y_1\cup Y_2\notin \mathcal{H}_j\setminus \mathcal{H}_{j-1}$.
Since $\mathcal{H}_j\setminus \mathcal{H}_{j-1}$ satisfies the condition $(*)$, at least one of $Y_1\cup Y_2$ and $Y_1\cap Y_2$ belongs to $\mathcal{H}_{j-1}$.
Since $Y_1,Y_2\in \mathcal{H}_j'$, we have $U_{j-1}\subseteq Y_1\cap Y_2\subseteq Y_1\cup Y_2\subseteq X$, which implies that there exists $Z\in \mathcal{H}_{j-1}$ with $U_{j-1}\subseteq Z\subseteq X$.
This contradicts the induction hypothesis.
\end{proof}

Take distinct $k$ elements $v_1,v_2,\ldots,v_k\in X$ that satisfy the condition in Claim \ref{clm:induction}.
Then, we have $Y\notin \mathcal{H}_k$ for all $Y\subseteq X$ with $U_k:=\{v_1,v_2,\ldots,v_k\}\subseteq Y$.
This implies that $Y\in \mathcal{F}$ for all $U_k\subseteq Y\subseteq X$.
Hence, setting $S=U_k$ satisfies the condition in Lemma \ref{lem:key}.
\end{proof}

\end{document}
